% ID: SEL-D0082-OB-E02
% ESTADO: cotejada. Segunda pasada visual y borrador compilado revisados.
% Fuente DOC: Selectividad/sel_2012_matematicas_aplicadas/4 MATEMÁT. APLIC. CC. SS/EXAMEN 4-ANDALUCÍA.doc
% SHA-256 original: 8446c5da6b4b7cfc0f48b89d1b2d052ae0c8a1c4b5bc0b86b7e63399ce8656f8
% Página convertida: 2; opción B; ejercicio 2.

Se estima que el beneficio de una empresa, en millones de euros, para los próximos 10 años viene dado por la función
\[
B(t)=\begin{cases}
at-t^2&\text{si }0\leq t\leq6\\
2t&\text{si }6<t\leq10
\end{cases},
\]
siendo $t$ el tiempo transcurrido en años.
\begin{enumerate}[label=\alph*)]
\item \textbf{(0.75 puntos)} Calcule el valor del parámetro $a$ para que $B$ sea una función continua.
\item \textbf{(1 punto)} Para $a=8$ represente su gráfica e indique en qué períodos de tiempo la función crecerá o decrecerá.
\item \textbf{(0.75 puntos)} Para $a=8$ indique en qué momento se obtiene el máximo beneficio en los primeros 6 años y a cuánto asciende su valor.
\end{enumerate}
